\section*{\texorpdfstring{\S\CurrentExerciseGroup}{第{\CurrentExerciseGroup}节}}
\phantomsection
\addcontentsline{toc}{section}{第{\CurrentExerciseGroup}节习题}

\begin{enumerate}
\def\labelenumi{\arabic{enumi})}
\item
  设 X 是一个局部紧空间，\(\Phi\) 是 \(\mathcal{M}(\mathrm{X};\mathbf{C})\) 上的一个滤子，使得 \(\mu\) 的支集沿 \(\Phi\) 退向无穷（No.~2，命题 7 之前）；试证：关于拟强拓扑（ \(\S1\) ，习题 8），\(\mu\) 关于 \(\Phi\) 趋于 0 。
\item
  试证：对于 \(\mathcal{M}^1 (\mathrm{X};\mathbf{C})\) 上的超强拓扑（§1，习题 15），紧支集测度所成之集在 \(\mathcal{M}^1 (\mathrm{X};\mathbf{C})\) 中稠密。
\item
  设 \(\mathbf{X}\) 是区间 {[}0,1{]} ，含于 \(\mathbf{R}\) 。试证：在 Banach 空间 \(\mathcal{M}(\mathbf{X};\mathbf{C})\) 中，Lebesgue 测度到每个离散测度的距离 \(\geqslant 1\) 。
\item
  设 A 是一个不可数集。在集 \(\mathfrak{P}(A)\) 中，考虑 M 与 N 之间的关系 \(\langle M \cap \complement N\) 与 \(N \cap \complement M\) 均可数 \(\rangle\) ；试证这是一个等价关系，且与包含关系 \(M \subset M'\) 弱相容（S, III, §1, 习题 2）。设 B 是 \(\mathfrak{P}(A)\) 关于这个等价关系的商集；试证：在 B 上，由包含关系通过过渡到商得到的关系是一个序关系，对于它 B 是一个 Boole 代数（loc. cit, 习题 17）。已知（GT, II, §4, 习题 12）存在 B 到一个 Boole 代数上的序结构同构，该 Boole 代数由一个全不连通紧空间 X 中既开又闭的子集组成。试证：若 U 是 X 中的非空开集，则存在由 U 的两两不相交的非空开子集 \(V_{\alpha}\) 组成的不可数族（利用下述事实：一个不可数集容许一个由不可数集组成的不可数分割）。由此推出：X 上的每个测度都具有无处稠密的支集。
\item
  设 \(X\) 是紧空间。称无穷序列 \((x_{n})_{n\in \mathbf{N}}\) （由 \(X\) 的点组成）
\end{enumerate}

具有极限分布，如果有限支集测度的序列 \(\left(\sum_{i=0}^{n-1} \varepsilon_{x_i}\right)/n\) 淡收敛到一个极限 \(\mu\) （在 \(\mathcal{M}_{+}(\mathrm{X})\) 中）；也称序列 \((x_{n})\) 关于测度 \(\mu\) 等分布。若 \(\mathrm{T}\) 是 Banach 空间 \(\mathcal{C}(\mathrm{X};\mathbf{C})\) 中的全子集，试证：为使序列 \((x_{n})\) 具有极限分布，必须且只须序列 \(\left(\left(\sum_{i=0}^{n-1} f(x_i)\right)/n\right)\) 在 \(\mathbf{C}\) 中收敛到一个极限（对每个函数 \(f \in \mathrm{T}\) ）。特别地，试证：若 \(X\) 是可度量的紧空间，则对每个序列 \((x_{n})\) （由 \(X\) 的点组成），存在具有极限分布的子序列 \((x_{n_k})\) 。

若 X 是 R 的区间 {[}0,1{]} 且 \(\mu\) 是 X 上的 Lebesgue 测度，试证：为使 X 的点的序列 \((x_{n})\) 关于 \(\mu\) 等分布，必须且只须它在 GT, VII, §1, 习题 14 的意义下 mod 1 等分布。为此，必须且只须对每个整数 \(m \in Z\) ，序列

\(\left(\left(\sum_{k=0}^{n-1} e^{2i\pi mx_k}\right)/n\right)\) 对 \(m \neq 0\) 趋于 0 。由此特别重新推出：序列 \(x_{n} = n\theta - [n\theta]\) （ \(n \in \mathbf{N}\) ）对每个无理数 \(\theta\) mod 1 等分布（GT, VII, §1, 习题 14）。当 \(\theta\) 为有理数时，其极限分布是什么？并参见 Ch. IV, §5, No.~12。

\setcounter{exercisegroup}{3}
\edef\CurrentExerciseGroup{\arabic{exercisegroup}}
